An der Erdoberfläche steht das Wasser unter dem Luftdruck, etwa dem Normdruck $p_\mathrm{n} = \pnvO$. Im luftleeren Raum fällt dieser Druck weg, $\Delta p = -p_\mathrm{n}$. Mit der Kompressibilität von Wasser $\chi = \chiWO$ (FoTaBe S.~188) und $\Delta V = -\chi\,V\,\Delta p$ (FoTaBe S.~164) dehnt sich das Wasser aus um
\begin{align}
 \Delta V &= \dVwF \\
   &= \chiW\cdot\Vwv\cdot\pnv \\
   &= \dVw \approx \result{\dVwmP{0}{2}}
\end{align}
Das Volumen nimmt also um rund \dVwcP{0}{2} zu, das ist der Bruchteil $\chi\,p_\mathrm{n} = \rVw$. (In Wirklichkeit würde das Wasser im Vakuum zu sieden beginnen und verdampfen; das ist hier nicht gefragt.)
